%%%%%%%%%%%%%%%%%%%%%%%%%%%%%%%%%%%%%%%%%%%%%%%%%%%%%%%%%%%%%%%%%%%%%%
% Lớp 8 | Chương I. Đa thức | Phép chia đa thức cho đơn thức
% Dạng 8D1?5-2: Chia đa thức cho đơn thức
%%%%%%%%%%%%%%%%%%%%%%%%%%%%%%%%%%%%%%%%%%%%%%%%%%%%%%%%%%%%%%%%%%%%%%


% nguồn: De_AI_Toan8_cuoi_hoc_ki_i_1 (nhập 21/11/2026)
\begin{bt}%[8D1V5-2]%[ma:626c49dd]
	Cho biểu thức $A = (x^{n+1}y^{n+2} - 2x^ny^{n+1}) : (x^ny^n)$ với $n \in \mathbb{N}$ và $x \ne 0$, $y \ne 0$. Rút gọn $A$ rồi tính giá trị của $A$ tại $x = 2$, $y = -1$.
	\loigiai{
		Ta có $A = (x^{n+1}y^{n+2} : x^ny^n) - (2x^ny^{n+1} : x^ny^n) = xy^2 - 2y$. \\
Thay $x = 2$, $y = -1$ vào $A$, ta được $A = 2 \cdot (-1)^2 - 2 \cdot (-1) = 2 + 2 = 4$.
	}
\end{bt}

% nguồn: De_AI_Toan8_cuoi_hoc_ki_i_1 (nhập 21/11/2026)
\begin{bt}%[8D1V5-2]%[ma:8489dd29]
	Thực hiện phép chia $\left[3(x - y)^4 + 2(x - y)^3 - 5(x - y)^2\right] : (y - x)^2$ (với $x \ne y$).
	\loigiai{
		Nhận xét rằng $(y - x)^2 = (x - y)^2$. \\
Đặt $t = x - y$ ($t \ne 0$), phép chia trở thành $(3t^4 + 2t^3 - 5t^2) : t^2 = 3t^2 + 2t - 5$. \\
Thay lại $t = x - y$, ta được kết quả là $3(x - y)^2 + 2(x - y) - 5$.
	}
\end{bt}

% nguồn: De_AI_Toan8_cuoi_hoc_ki_i_1 (nhập 21/11/2026)
\begin{bt}%[8D1V5-2]%[ma:2e41017f]
	Thực hiện phép chia $\left[4(x + 2y)^5 - 6(x + 2y)^3 + 2(x + 2y)^2\right] : \left[-2(x + 2y)^2\right]$ (với $x + 2y \ne 0$).
	\loigiai{
		Đặt $u = x + 2y$ ($u \ne 0$), biểu thức trở thành: \\
$(4u^5 - 6u^3 + 2u^2) : (-2u^2) = -2u^3 + 3u - 1$. \\
Thay lại $u = x + 2y$, ta được kết quả là $-2(x + 2y)^3 + 3(x + 2y) - 1$.
	}
\end{bt}

% nguồn: De_AI_Toan8_cuoi_hoc_ki_i_1 (nhập 21/11/2026)
\begin{bt}%[8D1V5-2]%[ma:8f6c40b5]
	Cho biểu thức $A = (6x^4y^2 - 12x^3y^2 + 18x^2y^2) : (6x^2y^2)$ với $x \ne 0$, $y \ne 0$.
a) Rút gọn biểu thức $A$.
b) Chứng minh rằng $A > 0$ với mọi $x \ne 0$, $y \ne 0$.
	\loigiai{
		a) Với $x \ne 0$, $y \ne 0$, ta có $A = x^2 - 2x + 3$. \\
b) Ta viết $A = (x^2 - 2x + 1) + 2 = (x - 1)^2 + 2$. \\
Vì $(x - 1)^2 \ge 0$ với mọi $x$, nên $A = (x - 1)^2 + 2 \ge 2 > 0$ với mọi $x \ne 0$, $y \ne 0$.
	}
\end{bt}

% nguồn: De_AI_Toan8_cuoi_hoc_ki_i_1 (nhập 21/11/2026)
\begin{bt}%[8D1V5-2]%[ma:9e2f717b]
	Tìm $x \ne 0$, biết $(10x^3 - 15x^2) : (5x^2) + (6x^2 + 9x) : (3x) = 8$.
	\loigiai{
		Với $x \ne 0$, ta thực hiện các phép chia: \\
$(10x^3 - 15x^2) : (5x^2) = 2x - 3$; \\
$(6x^2 + 9x) : (3x) = 2x + 3$. \\
Phương trình trở thành $(2x - 3) + (2x + 3) = 8 \Rightarrow 4x = 8 \Rightarrow x = 2$ (thỏa mãn $x \ne 0$). Vậy $x = 2$.
	}
\end{bt}

% nguồn: De_AI_Toan8_cuoi_hoc_ki_i_1 (nhập 21/11/2026)
